\documentclass[11pt]{article}

\usepackage{amsmath}
\usepackage{amssymb}
\usepackage{amsthm}
\usepackage[margin=1in]{geometry}
\usepackage{hyperref}
% The catalogue entry is bilingual; xeCJK is loaded only so that the Chinese
% half of the verbatim quote can be typeset. Replace the font below with any
% CJK-capable font available on your system if PingFang SC is absent.
\usepackage{xeCJK}
\setCJKmainfont{PingFang SC}

\newtheorem{theorem}{Theorem}
\newtheorem{proposition}[theorem]{Proposition}
\newtheorem{corollary}[theorem]{Corollary}
\theoremstyle{remark}
\newtheorem{remark}[theorem]{Remark}

\title{A rule-3 disproof submission for conjecture \texttt{00000007121}\\
       \large The Hirsch bound and its alleged ``high-dimensional double configuration''}
\author{earthking11}
\date{2026-09-13}

\begin{document}
\maketitle

\begin{abstract}
Conjecture \texttt{00000007121} asserts, in a single sentence, that the
Hirsch conjecture's diameter bound \emph{holds} \emph{and} that a
\emph{counterexample} to that bound is the ``high-dimensional double
configuration''. Read literally, the statement conjoins a claim with its own
negation, and is therefore contradictory. Independently, the Hirsch bound is
known to be false: Santos (2012) constructed a $43$-dimensional polytope with
$86$ facets and diameter $44 > 86 - 43 = 43$. We give both grounds, formalise
the logical content, and state prominently the caveat that the entry may be a
catalogue-style garbling rather than a substantive mathematical assertion.
\end{abstract}

\section{The statement, quoted exactly}

The catalogue entry \texttt{conjectures/00000007121.md} reads, verbatim:

\begin{quote}
\textbf{English.} Definition: Convex polytopes as geometric bodies of finite
halfspace intersections. Conjecture: The Hirsch conjecture's diameter upper
bound of face count minus dimension holds, and the counterexample to the bound
is the high-dimensional double configuration. (high-dimensional counterexample
to the Hirsch bound)

\textbf{中文。} 定义：凸多面体：有限半空间交的几何体。猜想：Hirsch 猜想：图直径的上界为面数减维数且界的反例为高维双重构型。（Hirsch 反例高维双构型直径上界）
\end{quote}

Let $d$ denote the dimension, $f$ the number of facets, and $\operatorname{diam}$
the graph diameter of a convex polytope. The Hirsch conjecture asserts
\begin{equation}\label{eq:hirsch}
  \operatorname{diam} \;\le\; f - d .
\end{equation}
Write
\[
  A \;:=\; \bigl(\forall\, d\, f\, \operatorname{diam},\;
        \operatorname{diam} \le f - d\bigr),
  \qquad
  B \;:=\; \bigl(\exists\, d\, f\, \operatorname{diam},\;
        f - d < \operatorname{diam}\bigr).
\]
The entry states $A \land B$ (with the second conjunct further naming the
``high-dimensional double configuration'').

\section{Ground 1: the statement as written is self-contradictory}

Over $\mathbb{N}$ with truncating subtraction, $f - d < \operatorname{diam}$
is the negation of $\operatorname{diam} \le f - d$ (Lean's
\texttt{Nat.not\_le}). Hence
\[
  B \;\Longleftrightarrow\; \lnot A ,
\]
and the statement as written is $A \land \lnot A$. This is a statement-level
(logical) refutation, and it is fully formalisable in core Lean~4; see
\texttt{lean4/Main.lean}, where the headline result is
\[
  \texttt{hirsch\_as\_written\_false}:\;
  \lnot\bigl(A \land B\bigr),
\]
proved by a kernel-checkable \texttt{intro}/\texttt{exact} argument over
\texttt{Nat}, together with the pure propositional fact
$\lnot(P \land \lnot P)$ and its instantiation.

\begin{proposition}\label{prop:logic}
The two conjuncts of conjecture \texttt{00000007121} cannot both hold:
\[
  \lnot\bigl(A \land B\bigr).
\]
Moreover any witness to $B$ entails $\lnot A$.
\end{proposition}

\begin{proof}
For the first claim, suppose $A \land B$. From $B$ obtain
$d, f, \operatorname{diam}$ with $f - d < \operatorname{diam}$. From $A$
obtain $\operatorname{diam} \le f - d$ for the same triple. Since natural
subtraction yields $\lnot(\operatorname{diam} \le f - d)$ from
$f - d < \operatorname{diam}$, we reach a contradiction. The second claim is
the same argument without the appeal to $A$. Both are formalised in
\texttt{lean4/Main.lean} (\texttt{counterexample\_refutes\_bound},
\texttt{hirsch\_as\_written\_false},
\texttt{conjecture\_00000007121\_false}).
\end{proof}

\section{Ground 2: the Hirsch bound is genuinely false (literature)}

Independently of the phrasing issue, \eqref{eq:hirsch} is false. Santos
\cite{Santos2012} constructed a convex polytope of dimension $43$ with $86$
facets whose graph diameter is $44$. Since
\[
  86 - 43 = 43 \;<\; 44,
\]
this violates \eqref{eq:hirsch} and disproves the Hirsch conjecture. The
polynomial Hirsch conjecture remains open.

\begin{remark}
This is a published literature result. It cannot be reproduced locally, and
this submission does not attempt to do so. The file \texttt{reproduce.py}
checks only the arithmetic $86 - 43 = 43 < 44$ and states explicitly that
Santos's construction itself is not reproduced. The Lean development does
\emph{not} formalise Santos's polytope; it formalises only the logical
inconsistency of the statement as written.
\end{remark}

\section{Caveat: the entry may be a catalogue artefact}

We state this prominently and do not overstate the logical reading.

\begin{itemize}
  \item The sentence names a conjecture and an alleged disproof of it in one
    breath. This is characteristic of an automatically assembled catalogue
    entry that has concatenated a statement with a ``counterexample'' tag,
    rather than a deliberate mathematical assertion.
  \item ``High-dimensional double configuration'' is \textbf{not} the standard
    name for Santos's object. The standard terminology is \emph{Hirsch
    polytopes}, \emph{spindles}, and the \emph{Klee--Walkup} and
    \emph{Santos} constructions. The phrase appears to be a non-standard or
    garbled label.
  \item Consequently the ``self-contradiction'' of Ground 1 is arguably a
    \emph{phrasing artefact} rather than a substantive mathematical
    falsehood, even though it is a genuine consequence of the text as
    written.
\end{itemize}

We therefore present both grounds---the literal logical inconsistency and the
independently known mathematical falsity---and leave the maintainer to judge
which reading of the entry is intended.

\begin{thebibliography}{9}
\bibitem{Santos2012}
F.~Santos.
\newblock A counterexample to the Hirsch conjecture.
\newblock \emph{Annals of Mathematics} \textbf{176} (2012), no.~1, 383--412.
\end{thebibliography}

\end{document}
